% Part: many-valued-logic
% Chapter: syntax-and-semantics
% Section: semantic-notions

\documentclass[../../../include/open-logic-section]{subfiles}

\begin{document}

\olfileid{mvl}{syn}{sem}

\olsection{مفاهیم معناشناختی}

فرض کنید منطق چندارزشی~$\Log L$ با یک ماتریس داده شده باشد. آنگاه
می‌توانیم مفاهیم معناشناختی متعارف را برای~$\Log L$ تعریف کنیم.

\begin{defn} 
\begin{enumerate}
\item !!^a{formula} چون~$!A$ \emph{ارضاپذیر} است اگر برای
  $\pAssign{v}$ای داشته باشیم~$\pSat{v}{!A}$؛ و \emph{ارضاناپذیر} است
  اگر برای هیچ~$\pAssign{v}$ای~$\pSat{v}{!A}$ نباشد؛
\item !!^a{formula} چون~$!A$ یک \emph{همان‌گویی} است اگر برای همهٔ
  !!{valuation}های~$v$، $\pSat{v}{!A}$؛
\item اگر~$\Gamma$ مجموعه‌ای از~!!{formula}ها باشد،
  $\Gamma \Entails !A$ («$\Gamma$، $!A$ را نتیجه می‌دهد») اگر و تنها
  اگر برای هر~!!{valuation} چون~$\pAssign{v}$ که
  $\pSat{v}{\Gamma}$، داشته باشیم~$\pSat{v}{!A}$.
\item اگر~$\Gamma$ مجموعه‌ای از~!!{formula}ها باشد، $\Gamma$
  \emph{ارضاپذیر} است اگر~!!a{valuation} چون~$\pAssign{v}$ وجود داشته
  باشد که~$\pSat{v}{\Gamma}$؛ وگرنه~$\Gamma$ \emph{ارضاناپذیر} است.
\end{enumerate} 
\end{defn}

برخی واقعیت‌هایی که دربارهٔ این مفاهیم داریم همان واقعیت‌های حالت منطق
کلاسیک‌اند:

\begin{prop}
\ollabel{prop:semanticalfacts} 
\begin{enumerate} 
\item $!A$ همان‌گویی است اگر و تنها اگر
  $\emptyset \Entails !A$؛
\item اگر~$\Gamma$ ارضاپذیر باشد، هر زیرمجموعهٔ متناهی~$\Gamma$ نیز
  ارضاپذیر است؛
\item\ollabel{def:monotonicity}%
یکنوایی: اگر~$\Gamma \subseteq \Delta$ و~$\Gamma \Entails !A$، آنگاه
  $\Delta \Entails !A$ نیز؛
\item\ollabel{def:Cut}%
تعدی: اگر~$\Gamma \Entails !A$ و
  $\Delta \cup \{ !A\} \Entails !B$، آنگاه
  $\Gamma \cup \Delta \Entails !B$؛
\end{enumerate}
\end{prop}

\begin{proof}
تمرین.
\end{proof}

\begin{prob}
\olref[mvl][syn][sem]{prop:semanticalfacts} را ثابت کنید.
\end{prob}

در منطق کلاسیک می‌توانیم استلزام را به شرطی پیوند دهیم. برای نمونه،
\emph{وضع مقدم} معتبر است: اگر~$\Gamma \Entails !A$ و
$\Gamma \Entails !A \lif !B$، آنگاه~$\Gamma \Entails !B$. رابطهٔ مهم
دیگر میان~$\Entails$ و~$\lif$ در منطق کلاسیک قضیهٔ استنتاج معناشناختی
است: $\Gamma \Entails !A \lif !B$ اگر و تنها اگر
$\Gamma \cup \{!A\} \Entails !B$. این نتایج در منطق‌های چندارزشی
همواره برقرار \emph{نیستند}. برقراربودن آن‌ها به تابع صدق~$\tf{\lif}$
بستگی دارد.

\end{document}
